\section[Pontrjagin 对偶与 Cartier 对偶]{Pontrjagin 对偶与 Cartier 对偶（2 月 2 日）}

\srcpage{15}

前几天我们花了一些时间讨论 Pontrjagin 对偶。回顾一下：取一个局部紧阿贝尔群（locally compact abelian group）$G$，它的对偶是
\[
  \widehat{G} := \Hom_{\mathrm{Ab}}(G,U_1),
\]
它典范地具有局部紧阿贝尔群的结构。这个对偶有许多类似 Fourier 变换（Fourier transform）的性质。一侧的 $L^1$ 函数，对应于另一侧在无穷远处消失的连续函数；一侧的 $L^2$ 函数，对应到另一侧的 $L^2$ 函数。一侧的特征标与另一侧的点互换，平移与卷积（convolution）对应于乘法。$G$ 紧当且仅当 $\widehat{G}$ 离散（当然，反之亦然）。

例子：

\begin{itemize}
  \item $U_1$ 的对偶是 $\Z$：特征标 $U_1\to U_1$ 就是映射 $z\mapsto z^n$，以 $n\in\Z$ 为指标。相应地，$\Z$ 的对偶是 $U_1$。
  \item 设 $T$ 是紧环面（compact torus），定义为商 $\R^d/\Lambda = \Lambda\otimes_\Z U_1$，其中 $\Lambda\subset\Z^d$ 是满秩的格（lattice）。它有相伴的对偶格 $\Lambda^\vee\subset\Hom(\Z^n,\Z)$\yizhu{原文写作 $\Hom(\Z^n,\Z)$，按上下文应为 $\Hom(\Z^d,\Z)$；译文照原文保留。}{A（原文 p.15）}，而 $T$ 的对偶就是 $\Lambda^\vee$。
  \item $\R^d$ 是自对偶的——更确切地说，$\R^d$ 的对偶是 $(\R^d)^*$。更一般地，对有限维实向量空间 $V$，对偶是 $V^*$，即通常的线性对偶。
\end{itemize}

最后这个等同恢复了 Fourier 变换的理论。设 $x$ 是原侧 $\R$ 上的坐标，$t$ 是对偶侧 $\R$ 上的坐标；那么 $\R_x\times\R_t$ 上的典范特征标是 $\exp(2\pi ixt)$。把变换写成拉回—推前（pullback-pushforward）时，得到的就是通常的 Fourier 变换。

更一般地，若从 $V$ 出发并得到 Pontrjagin 对偶 $V^*$（即通常的线性对偶），典范配对是 $\exp(2\pi i\langle x,t\rangle)$，而 Fourier 变换是
\[
  f(x) = \int_V \widehat{f}(t)\, e^{2\pi i\langle x,t\rangle}\, dt.
\]

在 Fourier 变换中我们知道，对 $x$ 求导（在原侧）与乘以 $t$（在对偶侧）互换；这与群元素上的平移和乘法互换是同一回事。这是合理的：求导是无穷小变换。这是一个普遍的例子，说明一侧的群论如何与另一侧的几何互换。

一般地，若 $G$ 是阿贝尔 Lie 群，其 Lie 代数 $\mathfrak{g}$ 映到 $G$ 上向量场的 Lie 代数，而后者位于 $G$ 上微分算子代数之内。此外，「忘记结合代数结构、只看 Lie 代数」的函子与「由 Lie 代数 $\mathfrak{g}$ 自由构造泛包络代数 $U(\mathfrak{g})$」的函子互为伴随。由于 $\mathfrak{g}$ 阿贝尔，$U(\mathfrak{g})$ 就是通常的对称代数 $\Sym^*\mathfrak{g}$。由伴随性，$U(\mathfrak{g})$ 作为微分算子作用在 $C^\infty(G)$ 上，于是可以在 $\Spec(U(\mathfrak{g})) = \mathfrak{g}^*$ 上取层化。例如对 $G = \R_x$，$\mathfrak{g}^* = \R_t$ 可等同于 $\Spec \R[t]$，其中 $t$ 就是 $\frac{d}{dx}$——这解释了为什么求导变成乘以 $t$。因此在原侧，群论的每个方面都被同时对角化了，因为一切都是交换的。

\srcpage{16}

说到交换的东西，我们来谈谈量子力学——众所周知它是交换的。\footnote{这是个玩笑。}这里正是 Fourier 变换会出现的地方。对在流形 $M$ 上运动的粒子，其经典力学的相空间是余切丛 $T^*M$，局部坐标为 $q$（位置）与 $p$（动量），其中动量指向丛的方向。它们生成可观测量代数。

在量子力学中，我们把 $T^*M$ 换成只依赖一半变量（即只依赖位置）的 $L^2$ 函数。也就是说，Hilbert 空间是 $\mathcal{H} := L^2(M)$，可观测量包括 $M$ 上的微分算子，既有 $M$ 上的函数也有切向量，其中 $p_j = i\frac{d}{dq_j}$。若 $\xi\in TM$、$f\colon M\to\R$ 是函数，则有交换子关系
\[
  \xi f = f\xi + \hbar f'.
\]
这里 $\hbar$ 是 Planck 常数。\yizhu{原文把交换子写成 $\xi f = f\xi + \hbar f'$，按字面即 $[\xi,f] = \hbar f'$；译文照原文，未作改动。}{A（原文 p.16）}

我们想说「态看起来像可观测量的平方根」，这与我们在量子化时经典地扔掉了一半可观测量有关。Fourier 变换重新登场，是因为我们想把动量算子及其导数对角化。这在一般情况下意义不大；但若 $M$ 是阿贝尔 Lie 群——常常是 $\R^n$——我们就有交换向量场的自然基，以及一个互换 $L^2(\R^n_q)\cong L^2(\R^n_p)$ 的 Fourier 变换。这最终给出 $G$ 上的量子力学与 $\widehat{G}$ 上的量子力学之间一个有趣的非平凡同构。关键在于，这个同构在经典层面是看不到的。这可以说是 $1$ 维阿贝尔对偶或镜像对称的极简版本；在更高维、更不平凡的设定中，我们会听到这个故事的回声。

接下来几周，我们会讨论 Cartier 对偶——这是上述故事的代数版本——然后转到它在物理与数论中的主要表现：电磁对偶与类域论。一个要点是：这两者都可以看成某种 Fourier 变换。它们是 Langlands 纲领的阿贝尔情形。

\begin{center}$\bullet$\end{center}

Cartier 对偶讲的是代数群（algebraic group）。可以从几种角度理解它：像是 Lie 群里的东西，但换到代数几何里，例如簇范畴中的群对象；不过用点函子（functor of points）来思考代数群也很有帮助。簇 $X$ 定义了函子 $\mathrm{Alg}_k\to\mathrm{Set}$，由 $X(R) := \Hom(\Spec R, X)$ 给出，即 $X$ 的「$R$-点」集合。这里 $\mathrm{Alg}_k$ 表示交换 $k$-代数。

若 $X$ 是簇范畴中的群对象，这等价于说其点函子穿过 $\mathrm{Grp}\to\mathrm{Set}$ 分解。也就是说，$X$ 的 $R$-值点 $\{\Spec R\to X\}$ 构成一个群，且这些群随 $R$ 变化而相容。由 Yoneda 引理，$X$ 由其点函子决定。

\begin{example}
  考虑函子 $R\mapsto \mathrm{GL}_n(R)$。它是可表的，从而定义了一个代数群 $\mathrm{GL}_n$。对 $\mathrm{SL}_n$ 等其它熟悉的群，也可以玩同样的游戏。
\end{example}

与 Pontrjagin 对偶一样，我们更关心阿贝尔代数群。由 Yoneda 引理的魔法，这等价于要求点函子穿过 $\mathrm{Ab}\hookrightarrow\mathrm{Grp}\to\mathrm{Set}$ 分解。

\begin{example}
  \begin{enumerate}
    \item 加法群 $\mathbb{G}_a := \mathbb{A}^1$，它把代数 $R$ 送到 $\Spec R$ 上函数的阿贝尔群。
    \item 乘法群 $\mathbb{G}_m := \mathbb{A}^1\setminus 0$。这是 $\Aut(k)$。其点函子是 $\mathbb{G}_m(R) = R^\times$。这是我们心中 $U_1$ 的类比物：特征标就是到 $\mathbb{G}_m$ 的函数。
    \item 另一个稍怪的例子：取值在 $\Z$ 中的常值函子。
  \end{enumerate}
\end{example}

现在给定阿贝尔代数群 $G$，可以令 $\widehat{G} := \Hom(G,\mathbb{G}_m)$，这里指代数群的同态（如果你愿意，也可以说是 $\mathrm{Grp}$-值点函子的自然变换）。

\textbf{有用的例子。}

\begin{itemize}
  \item $\widehat{\Z} = \mathbb{G}_m$，正如拓扑情形：从 $\Z$ 出发的映射由 $1$ 的像决定，而 $1$ 可以到任何地方。
  \item 反之，$\widehat{\mathbb{G}}_m = \Hom(\mathbb{G}_m,\mathbb{G}_m)\cong \Z$。
  \item 若从 $\Z/n$ 出发，得到 $\mu_n$，即 $n$ 次单位根的群概形。
\end{itemize}

设 $\Lambda$ 是复向量空间 $V$ 中的格。那么 $\Lambda$ 的对偶是对偶环面 $T^\vee = \Lambda^\vee\otimes_\Z \mathbb{G}_m$，其中 $\Lambda^*\subset V^*$ 是对偶格。环面 $V/\Lambda$ 的 Cartier 对偶是 $\Lambda^\vee$。这是 Fourier 级数在代数几何设定下的类比。

\srcpage{17}

为避免正则性问题，暂且假定 $G$ 有限（有限平坦群概形，或 Artin 代数的谱……）。\footnote{这些技术性问题只在正特征下出现：此时可以有非正则的有限群概形，例如 Frobenius 映射 $\varphi\colon\mathbb{A}^1\to\mathbb{A}^1$ 的核，它没有闭点，却是长度为 $p$ 的非约化概形。另一个（反）例子是 $n$ 次单位根的群概形：在 $\C$ 上它正如你所料就是 $\Z/n$，但在正特征下可能长得不一样。}此时 $\widehat{G} = \Spec(\mathcal{O}(G)^*)$。代数 $\mathcal{O}(G)$ 带有来自 $G$ 的额外结构，即由乘法映射的拉回给出的余代数（coalgebra）结构，称为余乘法（comultiplication）：
\[
  \Delta := \mu^*\colon \mathcal{O}(G)\longrightarrow \mathcal{O}(G)\otimes_k \mathcal{O}(G).
\]
它在这里扮演卷积的角色。类似地，有由单位元的包含 $\id\colon \Spec k\to G$ 的拉回给出的余单位映射 $\mathcal{O}(G)\to k$。

于是 $\mathcal{O}(G)$ 同时有代数与余代数结构。它们相处得好吗？好！它们满足若干交换关系，这些数据合起来给出一个交换、余交换的 Hopf 代数。（$G$ 的阿贝尔假设正是余交换性所需要的。）事实上，$k$ 上有限阿贝尔群概形的理论，反变等价于交换、余交换、有限维 Hopf 代数的理论。

对偶在交换、余交换、有限维 Hopf 代数的范畴上定义了一个对合。把它过一遍 $\Spec$，就得到有限阿贝尔群概形的 Cartier 对偶。若要研究无限代数群，就得操心正则性，因为无限维向量空间的对偶要麻烦一些。

\begin{remark}
  为什么余乘法应当是卷积？设 $G$ 是有限阿贝尔群，周围没有任何代数几何。那么 $(\C[G], *)$ 可等同于 $G$ 上测度的代数，而卷积等同于乘法映射下测度的推前。
\end{remark}

现在看看 Fourier 变换在这里会怎样。与拓扑情形一样，我们期待 $\mathbb{G}_a$ 映到 $\mathbb{G}_a$；但 $\mathbb{G}_a$ 的 Cartier 对偶实际上是另一个东西——$\mathbb{G}_a$ 的形式完备化（formal completion），而它令人困惑地也写作 $\widehat{\mathbb{G}}_a$。眼下我用 $\widehat{\mathbb{G}}_a$ 表示 Cartier 对偶，用 $\mathbb{G}_a^\wedge$ 表示形式完备化。后者其实是原点所有无穷小邻域的并：
\[
  \mathbb{G}_a^\wedge = \bigcup_n \Spec k[t]/(t^n).
\]

假设 $k$ 特征为零，这样就不必为如下事情操心：把 $\mathbb{G}_a = \Spec k[x]$ 的特征标定义为
\[
  e^{xt} = \sum_n \frac{(xt)^n}{n!}.
\]
要在代数几何中允许这个，就必须要求该和为有限和——这正是我们转去形式完备化的原因：对偶坐标 $t$ 是幂零的，所以没问题。

更一般地，若 $V$ 是有限维 $k$-向量空间（它定义了一个阿贝尔代数群），其 Cartier 对偶是 $(V^*)^\wedge$，即对偶在 $0$ 附近的形式完备化；对偶仍由 $\exp(\langle x,t\rangle)$ 给出，这在形式意义下（即在原点的这个形式邻域里）说得通。

谱理论问的是 $G$-表示，而 $G$-表示等同于 $\mathcal{O}(G)$-余模，即 $\mathcal{O}(\widehat{G})$-模，也就是 $\widehat{G}$ 上的拟凝聚层。要字面成立，$G$ 必须有限；但一般情形也有相应版本。Cartier 对偶正是以这种方式给出表示的分解。

例如，$\Rep(\Z)$ 等同于 $k[z,z^{-1}]$-模：$\mathcal{O}(\mathbb{G}_m) = k[z,z^{-1}]$。于是单个矩阵（一个向量空间连同自同态）等同于直线上的拟凝聚层，即在 $\mathbb{A}^1 = \Spec k[z]$ 上。要求矩阵可逆，就是说 $0$ 不能是特征值；于是你其实是 $\mathbb{A}^1\setminus 0 = \mathbb{G}_m$ 上的层——正如所料，这就是 $\Z$ 的 Cartier 对偶。

反方向上，$\mathbb{G}_m$ 的表示等同于 $\Z$-分次向量空间，即离散阿贝尔群 $\Z$ 上的层。也就是说，若 $V$ 是 $\mathbb{G}_m$-表示，可写成
\[
  V = \bigoplus_{n\in\Z} V_n,
\]
其中 $V_n$ 是 $z\in\mathbb{G}_m$ 以 $z^n$ 作用的特征子空间，这给出 $\Z$-分次。Fourier 级数就是这样把乘法群的表示变成离散集上的层。

\srcpage{18}

\begin{example}
  另一个 Fourier 级数的代数例子（格与环面互换），搭起来要费点功夫，但很漂亮。这个例子在拓扑里。设 $M$ 是 $3$-流形，$G$ 是 $M$ 的 Picard 群，即 $M$ 上复线丛的同构类群。群结构是张量积。我们也可以等价地用主 $U_1$-丛，或用 $\pi_0(\mathrm{Map}(M,BU_1))$。$BU_1$ 的模型可以取空间 $\C P^\infty$ 或 Eilenberg--Mac Lane 空间 $K(\Z,2)$。

  由于 $BU_1 = K(\Z,2)$，到 $BU_1$ 的映射的同伦类自然地与 $H^2(M;\Z)$ 等同（作为阿贝尔群），该等同把线丛送到它的第一 Chern 类。眼下假定 $H^2(M;\Z)$ 自由；不过有挠时也能讲这个故事的某个版本。

  现在 $G := H^2(M;\Z)$ 是阿贝尔的，来算它的 Cartier 对偶。这是一个对偶环面
  \[
    \Lambda^\vee\otimes_\Z \mathbb{G}_m = \Hom(H_1(M),\mathbb{G}_m)\cong \Hom(\pi_1(M),\mathbb{G}_m).
  \]
  它可以等同于 $M$ 上平坦 $\C^\times$-丛的同构类群：给定一个丛，用平坦联络定义单值映射 $\pi_1(M)\to\C^\times$。于是我们得到 $\Pic(M)$ 与 $\mathrm{Loc}_{\C^\times}(M)$ 之间的对偶——这并不深刻，不过是把 $\Z^n$ 与 $(\C^\times)^n$ 互换罢了。把线丛换成 $T$-丛（$T$ 是环面）可以使之不依赖基的选取。记 $T^\vee$ 为对偶环面。

  同样的故事让我们把 $\Bun_T(M)$ 与 $\mathrm{Loc}_{T^\vee}(M)$ 等同，后者是平坦 $T^\vee$-丛的同构类群。选一组基，它们就是 $(\C^\times)^n$-丛。平坦联络照样由单值决定。这两者都可以看成范畴，但我们不用这一点，只把它们当作同构类的集合。

  所以这仍然只是格—环面的 Cartier 对偶，但看起来更有启发性：丛与局部系统更几何；事实上，研究 $3$-流形上的电磁对偶时就会遇到这个例子。若把这一想法用到算术拓扑（arithmetic topology），把数域看成 $3$-流形，你就得到一个 Cartier 对偶的实例，并且得到类域论的陈述！
\end{example}

这种阿贝尔对偶本质上是四维的故事，并在所有更低维中有化身。它是 $3$-流形故事的一部分。而在这个故事里，你所拥有的远不止同构类的集合，而是多得多的结构。
